\section{Lessons for a GPU port of the whole of ALF}
\label{sec:lessons}

\file{GPU\_PLAN\_V2.md} (2026-11-13) records what the first plan assumed and what the measurements showed. The lessons below are taken from it, with the evidence from this document.

\begin{enumerate}
\item \textbf{OpenMP offload is a good base, not a complete solution.} The bulk kernels and the data management in OpenMP compiled and ran on three vendors from one Fortran source. But each vendor needed a small native layer for the hot small operations: SYCL on Intel for asynchrony, HIP on AMD for launch cost. Plan for it behind the library boundary from the start (Section~\ref{sec:backends}).
\item \textbf{Keep one library boundary.} \texttt{alf\_la} with a host reference implementation made every backend testable against LAPACK, and kept vendor code out of the physics. New device operations should enter as \texttt{la\_*} routines with a host version.
\item \textbf{Measure each primitive before redesigning.} The assumed cost of a synchronisation on the Max~1550 (100 to 170~$\mu$s) was wrong (6 to 18~$\mu$s); the real cost was a 32~KB copy through a small staging buffer (380~$\mu$s, 22~$\mu$s with a larger buffer). ``One environment variable gained more than two code phases.''
\item \textbf{Algorithm before kernels.} The first plan put the delayed and submatrix updates into a late phase and expected the stabilisation to dominate. The rank-one update turned out to be the largest kernel, and the submatrix update in block mode became the production scheme, 2.4 to 5 times faster. The next large gain is again algorithmic (two-level submatrix update).
\item \textbf{Vendor solver quality varies; decide per routine by measurement.} cuSOLVER was adequate; oneMKL's QR was hybrid and rocSOLVER's launch-bound. Own QR kernels of a few hundred lines fixed both. The orthogonality of cuSOLVER's $Q$ needed a Newton--Schulz step.
\item \textbf{The model must be flattened.} Derived types with allocatable components and polymorphic classes do not go to the device. A flat mirror built once at setup, plus refusal of everything the mirror cannot express, kept the CPU code untouched. A port of all of ALF needs a mirror for general vertices (rank $r$, non-diagonal, continuous fields) and for the observables of each model, or a host fallback for observables.
\item \textbf{Refuse rather than partially offload.} A feature is either fully on the device or refused at setup with a reason. Partial residency multiplies transfers and makes the data flow hard to reason about.
\item \textbf{Keep the random numbers and the decisions on the host.} It costs a few bytes per block and makes same-seed comparison with the CPU possible, which is the strongest and cheapest verification there is.
\item \textbf{Make reductions deterministic.} Fixed-order reductions make two GPU runs bitwise identical, which turns every switch and refactoring into a bitwise test.
\item \textbf{Rounding-level tolerances and statistical gates are weak tests.} They were costly, noise-limited, and caught no port bug. Bitwise identity and same-seed trajectories decide; statistics are reported. Every check must first be shown to fail on a planted defect, and verification should cost a bounded share of the work (at most about 20\,\%).
\item \textbf{Single-chain speed and throughput are different targets.} One chain is latency-bound; a production layout with several chains per GPU is throughput-bound. A change that removes launches helps the first; only a change that removes device work helps the second. The production layout (ranks per GPU, MPS share, CCS, $\Nd$) is a measured per-platform setting.
\item \textbf{Toolchains are fragile.} On the Max~1550 only oneAPI 2025.1 worked; nvhpc 24.11 produced NaN; nvfortran miscompiled one routine at \texttt{-O2}; GCC's AMD runtime needed a patch. Keep per-platform pages with a tested toolchain, and recheck on every module change.
\item \textbf{Hardware fails.} A faulty GPU tile produced corrupted Green functions. A tight precision threshold with an automatic restart, and a short health test per tile before production, are part of a GPU production setup.
\end{enumerate}

For an upstream merge \file{GPU\_PLAN\_V2.md} (WP5) proposes an 11-patch series, each building and passing the CPU tests with \texttt{GPU=none} unchanged: \texttt{alf\_la} with the host backend; the configure switch; the CPU delayed/submatrix updates; \file{gpu\_state\_mod} with \texttt{gpu\_supported} and the production guard; wraps and stabilisation with CUDA; updates; observables and $\tau$ pass; the oneMKL backend; the HIP backend; time-dependent hopping; user documentation.
